\chapter{Local Fields, Covariance, and Microcausality}
\label{ch:local-fields-covariance-microcausality}
\ConventionsLorentzian


This chapter works in the Wightman field presentation introduced in
Chapter~\ref{ch:volume-iv-wightman-fields-reconstruction}: local fields are
operator-valued distributions on a common invariant dense domain, and every
point-field formula is shorthand for a continuous smeared matrix element.

\section{Field Multiplets As Operator-Valued Distributions}

\begin{definition}[Field multiplet as an operator-valued distribution]
\label{def:field-multiplet-operator-valued-distribution}
Let \(\mathcal D\subset\Hilb\) be a dense linear subspace. In this chapter
\(\operatorname{End}(\mathcal D)\) denotes the algebra of linear maps
\(\mathcal D\to\mathcal D\). A field multiplet is a collection
\(\widehat\Phi_A\), with \(A\) in a finite index set, such that each component
defines a continuous linear map
\[
  \widehat\Phi_A:C_c^\infty(\mathbb M^D)\to\operatorname{End}(\mathcal D),
  \qquad
  f\mapsto\widehat\Phi_A(f).
\]
Continuity is meant in the usual locally convex LF topology on
\(C_c^\infty(\mathbb M^D)\). This is the localization version of the
test-function convention in Chapter~\ref{ch:starting-data-local-qft}; when a
later spectral statement uses \(\mathcal S(\mathbb M^D)\), the additional
tempered-continuity hypothesis is part of that statement. The notation
\[
  \widehat\Phi_A(f)
  =
  \int_{\mathbb M^D}\dd^Dx\,f(x)\widehat\Phi_A(x)
\]
is the distributional notation for this linear map.
The symbol \(x\) in \(\widehat\Phi_A(x)\) is a spacetime label, not an
operator on \(\Hilb\).  The operator is obtained only after smearing the
distributional field against a test function.  This convention separates
spacetime localization data from the Hilbert-space operators that act on
states.
\end{definition}

\begin{definition}[Localization of a smeared field]
\label{def:localized-smeared-field}
The support of the smeared field \(\widehat\Phi_A(f)\) is controlled by
\(\operatorname{supp}f\). If \(\operatorname{supp}f\subset\mathcal O\), where
\(\mathcal O\subset\mathbb M^D\) is a bounded open region, the field is
localized in \(\mathcal O\). In the local-algebra framework of Chapter
\ref{ch:starting-data-local-qft}, this means that \(\widehat\Phi_A(f)\) is
affiliated with \(\Obs(\mathcal O)\) or belongs to the unbounded field algebra
whose bounded functions generate observables localized in \(\mathcal O\).
\end{definition}

\begin{definition}[Common invariant field domain]
\label{def:common-invariant-field-domain}
  A common invariant field domain for a family of fields is a dense subspace
  \(\mathcal D\subset\Hilb\) such that every smeared field
  \(\widehat\Phi_A(f)\) maps \(\mathcal D\) into itself, every finite product
  of smeared fields is defined on \(\mathcal D\), the vacuum vector belongs to
  \(\mathcal D\), and the Poincare representation maps \(\mathcal D\) into
  itself.  If adjoints are used, the adjoint field is required to preserve the
  same domain.
\end{definition}

The domain condition is what makes a product such as
\(\widehat\Phi_{A_1}(f_1)\cdots\widehat\Phi_{A_n}(f_n)\vac\) an actual
vector before any distributional limit is taken.  Matrix elements
\[
  (f_1,\ldots,f_n)
  \longmapsto
  \left\langle\psi,
  \widehat\Phi_{A_1}(f_1)\cdots\widehat\Phi_{A_n}(f_n)\chi
  \right\rangle
\]
are required to be continuous multilinear functionals of the test functions
for \(\psi,\chi\) in a stated dense domain.  This continuity is the analytic
content behind point-field notation.

\begin{remark}[Distributional matrix elements]
\label{rem:distributional-matrix-elements}
  Suppose the smeared products above are defined on a common invariant field
  domain and the displayed matrix elements are continuous multilinear
  functionals on \(C_c^\infty(\mathbb M^D)^n\).  The Schwartz kernel theorem
  identifies the completed projective tensor product
  \(C_c^\infty(\mathbb M^D)^{\widehat\otimes n}\) with
  \(C_c^\infty((\mathbb M^D)^n)\).  Thus the displayed matrix element is
  equivalently a distribution
  \(W_{\psi,\chi}^{A_1\cdots A_n}\in\mathcal D'((\mathbb M^D)^n)\) satisfying
  \[
    W_{\psi,\chi}^{A_1\cdots A_n}(f_1\otimes\cdots\otimes f_n)
    =
    \left\langle\psi,
    \widehat\Phi_{A_1}(f_1)\cdots\widehat\Phi_{A_n}(f_n)\chi
    \right\rangle .
  \]
  This is the functional-analytic meaning of the stated continuity
  hypothesis.
\end{remark}

\begin{remark}[Direction of point-field notation]
\label{rem:point-field-notation-direction}
The preceding remark fixes the direction of logic: point-dependent
expressions are compact notation for continuous smeared matrix elements.
Later analytic continuation, spectral representation, and LSZ extraction all
act on these distributions or on their Fourier transforms.
\end{remark}

\section{Poincare Covariance Of Fields}

Let
\[
  g=(a,\Lambda)\in \Poinc
\]
act on spacetime by \(gx=\Lambda x+a\). Let \(S\) be a finite-dimensional
representation of \(SO^+(1,D-1)\), or of its double cover for spinor fields,
acting on the field components:
\[
  S(\Lambda)_A{}^B.
\]

\begin{definition}[Poincare pullback on test functions]
\label{def:poincare-pullback-test-functions}
For \(f\in C_c^\infty(\mathbb M^D)\), define the transformed test function
\[
  f_g(y)=f(g^{-1}y)=f(\Lambda^{-1}(y-a)).
\]
The assignment \(f\mapsto f_g\) is the pullback by \(g^{-1}\). With the group
law \(h g:x\mapsto h(gx)\), it satisfies
\[
  (f_g)_h=f_{hg}.
\]
\end{definition}

\begin{definition}[Covariant field multiplet]
\label{def:covariant-field-multiplet}
A field multiplet \(\widehat\Phi_A\) is Poincare covariant with component
representation \(S\) when
\[
  U(g)^{-1}\widehat\Phi_A(f)U(g)
  =
  S(\Lambda)_A{}^B\,\widehat\Phi_B(f_g)
\]
for every \(g=(a,\Lambda)\), every test function \(f\), and every component
index \(A\), with summation over repeated component indices.
\end{definition}

In point-field notation this covariance law is written
\[
  U(a,\Lambda)^{-1}\widehat\Phi_A(x)U(a,\Lambda)
  =
  S(\Lambda)_A{}^B\,\widehat\Phi_B(\Lambda x+a).
\]
The smeared equation is the definition; the point-field equation is its compact
distributional form.  With
\(U(a,\mathbf 1)=\exp(\ii a_\mu P^\mu)=\exp(\ii a^\mu P_\mu)\), conjugation
with \(U^{-1}\) on the left is the convention in which the coordinate label is
translated by \(+a\).

For a scalar field \(S(\Lambda)=1\). The infinitesimal translation law is
\[
  [P_\mu,\widehat\Phi(f)]
  =
  -\ii\,\widehat\Phi(\partial_\mu f),
\]
equivalently
\[
  [P_\mu,\widehat\Phi(x)]
  =
  \ii\,\partial_\mu\widehat\Phi(x)
\]
in distributional notation.  The two signs are the same distributional
statement: for compactly supported \(f\),
\[
  \int \dd^Dx\,f(x)\,\partial_\mu\widehat\Phi(x)
  =
  -\int \dd^Dx\,(\partial_\mu f)(x)\widehat\Phi(x),
\]
so moving the derivative from the field distribution to the test function
produces the minus sign in the smeared commutator.  Lorentz generators act by
an orbital term on the spacetime argument and a spin term on the component
index.

\begin{proposition}[Infinitesimal translation covariance]
\label{prop:infinitesimal-translation-covariance}
Let \(\widehat\Phi\) be a scalar covariant field and assume that
\(a\mapsto U(a,\mathbf 1)\) is generated on the common domain by
\(U(a,\mathbf 1)=\exp(\ii a^\mu P_\mu)\). Then, as an identity of smeared
operators on \(\mathcal D\),
\[
  [P_\mu,\widehat\Phi(f)]
  =
  -\ii\,\widehat\Phi(\partial_\mu f).
\]
Equivalently, in distributional point notation,
\[
  [P_\mu,\widehat\Phi(x)]
  =
  \ii\,\partial_\mu\widehat\Phi(x).
\]
\end{proposition}

\begin{proof}
For a pure translation \(g=(a,\mathbf 1)\),
Definition~\ref{def:covariant-field-multiplet} gives
\[
  U(a)^{-1}\widehat\Phi(f)U(a)
  =
  \widehat\Phi(f_a),
  \qquad
  f_a(y)=f(y-a).
\]
Differentiate at \(a=0\) in the \(\mu\)-direction. The left side gives
\[
  \frac{d}{dt}\Big|_{t=0}
  e^{-\ii tP_\mu}\widehat\Phi(f)e^{\ii tP_\mu}
  =
  -\ii[P_\mu,\widehat\Phi(f)],
\]
while the right side gives
\[
  \frac{d}{dt}\Big|_{t=0}\widehat\Phi(f(\,\cdot-te_\mu))
  =
  -\widehat\Phi(\partial_\mu f).
\]
Equating the two derivatives gives the smeared formula. Pairing the
point-field identity with \(f\) and integrating by parts gives the
distributional version.
\end{proof}

\paragraph{Covariance of localization regions.}
If \(\operatorname{supp}f\subset\mathcal O\), then
\(\operatorname{supp}f_g\subset g\mathcal O\). Therefore the covariant
transform \(U(g)^{-1}\widehat\Phi_A(f)U(g)\) is localized in \(g\mathcal O\)
whenever \(\widehat\Phi_A(f)\) is localized in \(\mathcal O\).
The support identity \(\operatorname{supp}f_g=g\,\operatorname{supp}f\)
follows from \(f_g(y)=f(g^{-1}y)\) and continuity of the homeomorphism \(g\).
The covariance law then rewrites the transformed operator as a linear
combination of fields smeared with test functions supported in
\(g\mathcal O\).

\paragraph{Covariance of Wightman distributions.}
Assume that \(\widehat\Phi_A\) is a covariant field multiplet on the common
invariant domain of Definition~\ref{def:common-invariant-field-domain}, and
that the vacuum is Poincare invariant:
\[
  U(g)\vac=\vac,\qquad g\in\Poinc .
\]
Then the \(n\)-point Wightman distributions satisfy
\[
\begin{split}
  W_{A_1\cdots A_n}(f_1,\ldots,f_n)
  &=
  S(\Lambda)_{A_1}{}^{B_1}\cdots
  S(\Lambda)_{A_n}{}^{B_n}
\\
  &\quad \times
  W_{B_1\cdots B_n}\bigl((f_1)_g,\ldots,(f_n)_g\bigr),
  \qquad g=(a,\Lambda),
\end{split}
\]
where \((f_i)_g(y)=f_i(g^{-1}y)\).  In point-field notation this is
\[
  W_{A_1\cdots A_n}(x_1,\ldots,x_n)
  =
  S(\Lambda)_{A_1}{}^{B_1}\cdots S(\Lambda)_{A_n}{}^{B_n}
  W_{B_1\cdots B_n}(gx_1,\ldots,gx_n),
\]
with the same distributional meaning as the smeared formula.  The derivation
uses only the stated field covariance and vacuum invariance.
By vacuum invariance and unitarity,
\[
  \bra{\vac}\widehat\Phi_{A_1}(f_1)\cdots\widehat\Phi_{A_n}(f_n)\ket{\vac}
  =
  \bra{\vac}
  U(g)^{-1}\widehat\Phi_{A_1}(f_1)\cdots
  \widehat\Phi_{A_n}(f_n)U(g)
  \ket{\vac}.
\]
Since \(U(g)\mathcal D\subset\mathcal D\), conjugation by \(U(g)\) may be
applied to each factor on the common domain.  Using
Definition~\ref{def:covariant-field-multiplet} for every field gives
\[
  U(g)^{-1}\widehat\Phi_{A_i}(f_i)U(g)
  =
  S(\Lambda)_{A_i}{}^{B_i}\widehat\Phi_{B_i}((f_i)_g).
\]
Multiplying these identities in the given operator order yields the smeared
covariance formula.  The point-field formula is exactly the same identity
paired with tensor-product test functions.

\section{Spacelike Separation Of Supports}

\begin{definition}[Spacelike separated supports]
\label{def:spacelike-separated-supports}
Two subsets \(K_1,K_2\subset\mathbb M^D\) are spacelike separated when
\[
  (x-y)^2>0
  \qquad
  \text{for every }x\in K_1,\ y\in K_2 .
\]
Thus two test functions \(f,g\in C_c^\infty(\mathbb M^D)\) have spacelike
separated supports when \(K_1=\operatorname{supp}f\) and
\(K_2=\operatorname{supp}g\) satisfy this condition.
\end{definition}

\paragraph{Poincare invariance of spacelike support separation.}
Let \(g=(a,\Lambda)\in\Poinc\). If \(K_1,K_2\subset\mathbb M^D\) are
spacelike separated, then \(gK_1\) and \(gK_2\) are spacelike separated:
for \(x\in K_1\) and \(y\in K_2\),
\[
  (gx-gy)^2=(\Lambda(x-y))^2=(x-y)^2>0,
\]
because translations cancel in differences and \(\Lambda\) preserves the
Minkowski quadratic form.  For test functions,
\(\operatorname{supp}(f_g)=g\,\operatorname{supp}f\):
indeed \(f_g(y)=f(g^{-1}y)\) can be nonzero only when
\(g^{-1}y\in\operatorname{supp}f\), equivalently
\(y\in g\,\operatorname{supp}f\), and the reverse inclusion follows after
taking closures.  Therefore Poincare transformation preserves spacelike
separation of supports.

\begin{center}
\begin{tikzpicture}[scale=0.95,>=Stealth]
  \draw[->,line width=0.7pt] (0,-2.5) -- (0,2.65) node[above] {$x^0$};
  \draw[->,line width=0.7pt] (-3.2,0) -- (3.45,0) node[right] {$x^1$};
  \draw[green!45!black,line width=1.0pt] (0,0) -- (-1.8,2.25);
  \draw[green!45!black,line width=1.0pt] (0,0) -- (1.8,2.25);
  \draw[green!45!black,line width=1.0pt] (0,0) -- (-1.8,-2.25);
  \draw[green!45!black,line width=1.0pt] (0,0) -- (1.8,-2.25);
  \node[green!45!black] at (2.55,2.05) {light cone};
  \fill[blue!18] (-2.1,0.45) ellipse (0.52 and 0.34);
  \draw[blue!70!black,line width=0.9pt] (-2.1,0.45) ellipse (0.52 and 0.34);
  \node[blue!70!black,anchor=east] at (-2.72,0.45) {$\operatorname{supp}f$};
  \fill[red!18] (2.1,0.45) ellipse (0.52 and 0.34);
  \draw[red!70!black,line width=0.9pt] (2.1,0.45) ellipse (0.52 and 0.34);
  \node[red!70!black,anchor=west] at (2.72,0.45) {$\operatorname{supp}g$};
  \draw[<->,gray!70,line width=0.8pt] (-1.55,0.45) -- (1.55,0.45);
  \node[gray!70!black,fill=white,inner sep=1.2pt] at (0,1.18)
    {spacelike separation};
\end{tikzpicture}
\end{center}

The figure displays one spatial dimension. The condition is the set-theoretic
one written above: every pair of points, one from each support, has spacelike
separation.

\section{Microcausality}

\begin{definition}[Homogeneous parity and graded commutator]
\label{def:homogeneous-parity-graded-commutator}
Suppose the field multiplets carry a fermion parity \(|\widehat\Phi_A|\in
\{0,1\}\). For homogeneous operators \(X,Y\), define the graded commutator by
\[
  [X,Y]_{\mathrm{gr}}
  =
  XY-(-1)^{|X||Y|}YX.
\]
For bosonic fields both parities are zero, so this is the ordinary
commutator.
\end{definition}

\begin{definition}[Microcausality for smeared fields]
\label{def:microcausality-smeared-fields}
The field multiplets satisfy microcausality on \(\mathcal D\) when
\[
  [\widehat\Phi_A(f),\widehat\Psi_B(g)]_{\mathrm{gr}}\psi=0
\]
for every \(\psi\in\mathcal D\), whenever
\(\operatorname{supp}f\) and \(\operatorname{supp}g\) are spacelike separated.
\end{definition}

In point-field notation the bosonic case is written
\[
  [\widehat\Phi_A(x),\widehat\Psi_B(y)]=0
  \qquad
  (x-y)^2>0.
\]
The smeared statement is the definition. It is the local compatibility
condition for spacelike separated field insertions.

\paragraph{Adjacent spacelike exchange in Wightman functions.}
Assume microcausality in the sense of
Definition~\ref{def:microcausality-smeared-fields}.  Let
\(\widehat\Phi_A\) and \(\widehat\Psi_B\) be homogeneous fields of parities
\(|A|\) and \(|B|\).  If \(f\) and \(g\) have spacelike separated supports,
then for any fields placed to the left and right whose products are defined
on the common domain,
\[
\begin{split}
  &\bra{\vac}
  \widehat X\,
  \widehat\Phi_A(f)\widehat\Psi_B(g)\,
  \widehat Y
  \ket{\vac}
\\
  &\qquad =
  (-1)^{|A||B|}
  \bra{\vac}
  \widehat X\,
  \widehat\Psi_B(g)\widehat\Phi_A(f)\,
  \widehat Y
  \ket{\vac}.
\end{split}
\]
Equivalently, adjacent entries of a Wightman distribution may be exchanged
across a spacelike support separation with the Koszul sign
\((-1)^{|A||B|}\).

Microcausality gives, on \(\mathcal D\),
\[
  \widehat\Phi_A(f)\widehat\Psi_B(g)
  -
  (-1)^{|A||B|}
  \widehat\Psi_B(g)\widehat\Phi_A(f)=0 .
\]
Let \(\widehat Y\ket{\vac}\in\mathcal D\); this is part of the common-domain
assumption for finite products of smeared fields.  Applying the displayed
operator identity to this vector and then pairing with
\(\bra{\vac}\widehat X\) gives the stated matrix-element identity.  Since
Wightman distributions are the continuous multilinear extensions of these
smeared matrix elements, the same exchange relation is their distributional
support-separated boundary value.

\section{Field Coordinates And Local Algebras}

\begin{definition}[Local polynomial field algebra on a common domain]
\label{def:local-polynomial-field-algebra}
For a bounded open region \(\mathcal O\), let
\[
  \mathcal P(\mathcal O)
\]
denote the polynomial \(*\)-algebra on \(\mathcal D\) generated by finite sums
and products of smeared fields \(\widehat\Phi_A(f)\) with
\(\operatorname{supp}f\subset\mathcal O\), together with their adjoints when
the adjoints preserve the chosen domain. This algebra records local unbounded
field coordinates.
\end{definition}

\paragraph{Isotony and graded locality for polynomial field algebras.}
If \(\mathcal O_1\subset\mathcal O_2\), then
\(\mathcal P(\mathcal O_1)\subset\mathcal P(\mathcal O_2)\). If
\(\mathcal O_1\) and \(\mathcal O_2\) are spacelike separated and the fields
satisfy Definition~\ref{def:microcausality-smeared-fields}, then homogeneous
monomials \(X\in\mathcal P(\mathcal O_1)\) and
\(Y\in\mathcal P(\mathcal O_2)\) satisfy
\[
  XY=(-1)^{|X||Y|}YX
\]
on the common invariant domain, after expanding \(X\) and \(Y\) into ordered
products of homogeneous smeared fields.
The isotony statement follows directly from support inclusion: every test
function supported in \(\mathcal O_1\) is also supported in \(\mathcal O_2\).
For locality, every generating smeared field from
\(\mathcal P(\mathcal O_1)\) has spacelike separated support from every
generating smeared field from \(\mathcal P(\mathcal O_2)\). Microcausality
therefore permits each generator in \(X\) to be moved past each generator in
\(Y\). Each exchange of homogeneous factors contributes the sign
\((-1)^{|A||B|}\). Multiplying these signs over all crossings gives the total
Koszul sign \((-1)^{|X||Y|}\), where \(|X|\) and \(|Y|\) are the sums of the
parities of their constituent fields modulo \(2\).

The observable algebra \(\Obs(\mathcal O)\subset\mathcal B(\Hilb)\) is the
bounded local algebra used in the opening framework. The relation between
\(\mathcal P(\mathcal O)\) and \(\Obs(\mathcal O)\) depends on the analytic
framework. For free scalar fields, Weyl operators are bounded functions of
smeared fields and generate a local algebra. In Wightman language one works
directly with operator-valued distributions on a common domain. In algebraic
language one may take the bounded local net as primary and use fields as
coordinates, charged fields, or affiliated unbounded operators when available.

The monograph will use whichever of these frameworks is required by the
construction at hand, with the framework declared before theorem-level claims
are made.

\section{Wightman Functions}

\begin{definition}[Wightman distribution of a field multiplet]
\label{def:wightman-distribution-field-multiplet}
Let \(\widehat\Phi_{A_1},\ldots,\widehat\Phi_{A_n}\) be fields whose products
on the vacuum vector are defined on the common domain. The \(n\)-point
Wightman distribution is the multilinear functional
\[
  W_{A_1\cdots A_n}(f_1,\ldots,f_n)
  =
  \bra{\vac}
  \widehat\Phi_{A_1}(f_1)\cdots
  \widehat\Phi_{A_n}(f_n)
  \ket{\vac}.
\]
Point notation writes this distribution as
\[
  W_{A_1\cdots A_n}(x_1,\ldots,x_n)
  =
  \bra{\vac}
  \widehat\Phi_{A_1}(x_1)\cdots
  \widehat\Phi_{A_n}(x_n)
  \ket{\vac}.
\]
\end{definition}

\paragraph{Field axioms inherited by Wightman distributions.}
Assume the common-domain and continuity hypotheses of
Remark~\ref{rem:distributional-matrix-elements}.  If the vacuum is
Poincare invariant and the field multiplet is covariant, inserting
\(U(g)U(g)^{-1}\) between neighboring smeared factors and using vacuum
invariance gives the covariance law displayed in the covariance paragraph
above.  If the
field multiplet satisfies microcausality, applying the vanishing graded
commutator to adjacent factors with spacelike separated supports gives the
graded exchange relation displayed above.  The continuity of the smeared
matrix elements then extends these identities from tensor products of test
functions to the corresponding Wightman distributions.

Poincare covariance of the fields gives covariance of the Wightman
distributions. Microcausality gives graded exchange relations when adjacent
insertions have spacelike separated supports. The spectrum condition will give
the analytic and spectral properties used in the K\"all\'en--Lehmann
representation.

\section{Position In The Construction}

The data introduced in this chapter are local field data: test functions,
smeared operators, covariance, and microcausality. They prepare the next stage,
where Lagrangian and path-integral methods are introduced first as
constructions of Green functions. Spectral representation will then connect
two-point functions to particle content before scattering theory is defined.
